\capitulocapa{images/cap-4.png}{Por que Context API nem sempre resolve?}
\label{chap:context-api}

\section{O Provider que virou gaveta}

\begin{historia}
A migração do tema foi bem recebida. Sumiram várias propriedades, o Header ficou
menor e o Perfil passou a acessar a preferência sem conhecer o Layout.

Então veio a ideia que parecia economizar tempo:

``Já temos um contexto global. Coloca nele também os filtros do Dashboard, o
usuário selecionado, o Drawer e as notificações.''

Algumas semanas depois, o arquivo tinha centenas de linhas. Uma alteração no
campo de busca fazia componentes distantes renderizarem novamente. Ninguém sabia
quais ações pertenciam a qual parte do estado.

Context continuava funcionando. O problema era o trabalho que passamos a exigir
dele.
\end{historia}

Context API resolveu o \emph{como distribuir} o tema. Ela não decidiu quais
dados deveriam ser globais, como organizar ações, com que frequência o valor
mudaria ou quais consumidores deveriam reagir a cada mudança.

Essa distinção evita dois extremos: rejeitar Context por completo ou transformá-la
automaticamente no store da aplicação.

\section{O problema: tudo cabe no mesmo contexto}

Tecnicamente, nada impede um contrato como este:

\begin{lstlisting}[caption={Um contexto que começou a acumular responsabilidades.}]
type AppContextValue = {
  theme: Theme
  toggleTheme: () => void
  search: string
  setSearch: (value: string) => void
  status: UserStatus
  setStatus: (status: UserStatus) => void
  isDrawerOpen: boolean
  openDrawer: () => void
  closeDrawer: () => void
}
\end{lstlisting}

O nome \codigo{AppContext} já é um sinal. Ele não descreve um domínio ou uma
responsabilidade; descreve apenas o fato de estar na aplicação.

Tema, busca e Drawer possuem motivos diferentes para mudar. Também possuem
consumidores e ciclos de vida diferentes. Colocá-los no mesmo objeto cria uma
fronteira artificial: qualquer mudança produz um novo valor para o mesmo
Provider.

\section{Como normalmente tentamos corrigir}

Quando aparecem renderizações inesperadas, a primeira tentativa costuma ser
envolver o valor em \codigo{useMemo}:

\begin{lstlisting}
const value = useMemo(
  () => ({
    theme,
    toggleTheme,
    search,
    setSearch,
    isDrawerOpen,
    openDrawer,
    closeDrawer,
  }),
  [theme, search, isDrawerOpen],
)
\end{lstlisting}

Isso evita criar outro objeto quando nenhuma dependência mudou. Mas, quando a
pessoa digita uma letra na busca, \codigo{search} muda, o objeto muda e os
consumidores daquele contexto recebem um novo valor.

Memoização não cria seleção por campo. Ela não permite que um consumidor assine
somente \codigo{theme} dentro de um objeto cujo \codigo{search} muda a cada
tecla.

Outra tentativa é envolver todos os componentes em \codigo{memo}. Um consumidor
de contexto, porém, precisa reagir quando o contexto que ele lê muda. A
memoização de propriedades não transforma o contexto numa assinatura granular.

\begin{decisao}
Antes de adicionar memoização, verifique se a fronteira está correta. Otimizar um
contexto que mistura responsabilidades preserva a causa estrutural e acrescenta
mais código para o time entender.
\end{decisao}

\section{Por que as renderizações se espalham?}

O Provider disponibiliza um valor. Quando esse valor muda, React precisa
notificar os componentes que consomem aquele contexto. O consumidor pode usar
apenas um campo, mas sua inscrição é no contexto, não naquele campo isolado.

Considere dois componentes:

\begin{lstlisting}
function ThemeBadge() {
  const { theme } = useAppContext()
  return <span>{theme}</span>
}

function SearchField() {
  const { search, setSearch } = useAppContext()
  return <input value={search} onChange={(event) => {
    setSearch(event.target.value)
  }} />
}
\end{lstlisting}

Ao digitar, o valor do Provider muda. \codigo{ThemeBadge} volta a executar mesmo
que \codigo{theme} continue igual. Em uma aplicação pequena, isso pode ser
irrelevante. Em telas grandes, com mudanças frequentes e muitos consumidores, o
custo começa a aparecer.

\ilustracaopendente{ch04-renderizacao-context}{Uma alteração em search cria novo valor no AppContext e alcança ThemeBadge, Sidebar e SearchField, destacando consumidores sem mudança útil.}

O objetivo não é perseguir zero renderizações. Renderizar é parte do React. O
problema é acoplar regiões independentes a uma frequência de mudança que não
pertence a elas.

\section{Mistura de responsabilidades}

Desempenho é apenas uma parte do problema. Um contexto grande também mistura
decisões:

\begin{itemize}
  \item quem pode alterar cada campo;
  \item quais valores devem persistir;
  \item o que deve ser reiniciado ao trocar de rota;
  \item quais ações representam regras do domínio;
  \item como testar cada comportamento isoladamente.
\end{itemize}

O Provider passa a conhecer detalhes de interface, preferências, filtros e
sessão. Um teste do tema precisa montar dependências da busca. Uma alteração no
Drawer modifica o mesmo contrato usado pelo Perfil.

Context não causou essa mistura. Nós escolhemos a fronteira errada porque o
mecanismo aceitava qualquer valor.

\section{Uma solução melhor: contextos pequenos}

O tema continua sendo um bom candidato para Context:

\begin{itemize}
  \item é transversal;
  \item muda pouco;
  \item possui contrato pequeno;
  \item seus consumidores estão dentro da mesma área da aplicação;
  \item não exige seleção complexa.
\end{itemize}

Portanto, mantemos \codigo{ThemeProvider} pequeno. Os filtros do Dashboard podem
permanecer próximos da página que os utiliza. O Drawer pode ganhar outra
fronteira quando houver consumidores suficientes. Não precisamos escolher uma
única ferramenta para todos os estados.

\begin{lstlisting}[caption={Um Provider com uma única responsabilidade.}]
type ThemeContextValue = {
  theme: Theme
  toggleTheme: () => void
}

function ThemeProvider({ children }: ThemeProviderProps) {
  const [theme, setTheme] = useState<Theme>('light')

  function toggleTheme() {
    setTheme((current) => (current === 'light' ? 'dark' : 'light'))
  }

  return (
    <ThemeContext value={{ theme, toggleTheme }}>
      {children}
    </ThemeContext>
  )
}
\end{lstlisting}

Separar contextos não significa criar um Provider para cada booleano. A
fronteira deve representar uma responsabilidade coesa. Tema e ação de alternar
mudam pelo mesmo motivo; busca e Drawer não.

\section{Estado local ainda participa da arquitetura}

Os filtros do Dashboard não precisam ser globais apenas porque mais de um
componente da página os utiliza. Podemos mantê-los na própria página e
distribuí-los para os controles de busca e de status.

\begin{lstlisting}
function DashboardPage() {
  const [search, setSearch] = useState('')
  const [status, setStatus] = useState<UserStatus>('all')

  return (
    <DashboardFilters
      search={search}
      status={status}
      onSearchChange={setSearch}
      onStatusChange={setStatus}
    />
  )
}
\end{lstlisting}

Essa solução mantém a frequência de digitação dentro da página. O Header e a
Sidebar não participam. Se outra rota precisar dos mesmos filtros no futuro,
reavaliaremos a fronteira com um requisito concreto.

\section{Quando Context ainda faz sentido}

Context costuma funcionar bem quando:

\begin{itemize}
  \item a informação é naturalmente transversal;
  \item o contrato é pequeno e estável;
  \item as mudanças são pouco frequentes;
  \item a maioria dos consumidores se interessa pelo valor completo;
  \item o Provider possui responsabilidade clara.
\end{itemize}

Tema, idioma e dependências de infraestrutura são exemplos comuns. Sessão pode
ser adequada, desde que credenciais, dados remotos e preferências não sejam
misturados sem critério.

Context começa a incomodar quando precisamos de seletores, ações organizadas,
leitura fora do React, middlewares ou muitas atualizações independentes. Esses
requisitos não tornam Context ruim; indicam que precisamos de uma ferramenta
com outro modelo.

\section{Exercício: faça o inventário}

\begin{tarefa}
O time sugeriu colocar tema, filtros do Dashboard, Drawer de notificações,
usuário da sessão e rascunho do Perfil em um único \codigo{AppContext}. Antes de
implementar, classifique cada item e decida onde ele deve viver por enquanto.
\end{tarefa}

Para cada estado, registre:

\begin{enumerate}
  \item quem lê e quem altera;
  \item com que frequência muda;
  \item por quanto tempo deve existir;
  \item se pertence a uma tela, à interface inteira ou ao servidor;
  \item se Context oferece uma fronteira coesa para ele.
\end{enumerate}

\espacoresposta{10}

Não procure uma resposta única chamada ``estado global''. O exercício está
correto quando cada decisão possui um motivo que o restante do time consegue
questionar e revisar.

\begin{resumo}
\begin{itemize}
  \item Context distribui um valor; ela não organiza automaticamente o estado.
  \item Consumidores reagem quando o valor do contexto muda, mesmo quando usam
    apenas uma parcela que permaneceu igual.
  \item Memoização não corrige uma fronteira que mistura responsabilidades.
  \item Contextos pequenos reduzem acoplamento e tornam contratos mais claros.
  \item Estado local continua sendo a primeira opção para demandas restritas a
    uma tela.
  \item Context é uma boa escolha para valores transversais, pequenos e pouco
    mutáveis.
\end{itemize}
\end{resumo}

No próximo capítulo, separaremos estado do cliente e estado do servidor. Essa
distinção evitará colocar no contexto outro tipo de dado que já possui seus
próprios problemas de sincronização.
